% !TEX root = ../../supporting_information.tex

\section{Hierarchical Uncertainty-Quantification Details}
\label{sec:si-hierarchical-uq}

\subsection{Within-window joint permeability states}

For each window, component-wise empirical ranks are computed locally within its PREDICT library. The joint rank score for realization $m$ is

\begin{equation}
  r_{\mathrm{joint}}^{(m)} =
  \frac{r_{xx}^{(m)} + r_{yy}^{(m)} + r_{zz}^{(m)}}{3}.
  \label{eq:si-joint-rank-score}
\end{equation}

The score orders clusters from relatively low to relatively high joint permeability; clustering itself is performed in physical three-dimensional $\log_{10}(k)$ space. \todo{Add the final clustering settings, admissibility rules, silhouette threshold, state-library fraction, medoid definition, and local/state-wide pool construction.}

\subsection{Window similarity groups}

For windows with empirical libraries $A$ and $B$, the multivariate energy distance is

\begin{equation}
  E(A,B) = 2\,\mathbb{E}\lVert X-Y\rVert
  - \mathbb{E}\lVert X-X'\rVert
  - \mathbb{E}\lVert Y-Y'\rVert,
  \label{eq:si-energy-distance}
\end{equation}

where $X,X'\sim A$ and $Y,Y'\sim B$ are independent samples in physical $\log_{10}(k)$ space. \todo{Define the within-window normalization, similarity threshold, all-pairs grouping rule, tie-breaking criteria, and optional bootstrap stability assessment exactly as implemented.}

\subsection{Multiple-window permeability cases}

\todo{Tabulate the ten case types: two independent cases, four fault-wide low/high cases, and four grouped low/high cases. Explain that similarity grouping supports candidate coupling assumptions but does not infer true cross-window correlation. State how singleton groups remain independently sampled.}

% Place supporting workflow and sensitivity figures in supplement/figures and
% use Figure S1, Figure S2, ... automatically through the numbering above.
